% ---- score2 standard preamble (from: python3 score2_tools.py preamble) ----
\documentclass[10pt]{extarticle}
\usepackage[a4paper,margin=1.5cm]{geometry}
\usepackage[T1]{fontenc}
\usepackage{amsmath,amssymb}
\usepackage{graphicx}
\usepackage[export]{adjustbox}
\usepackage{multicol}
\usepackage{enumitem}
\usepackage{needspace}
\usepackage{booktabs,array,tabularx}
\setlength{\columnsep}{0.7cm}
\setlength{\parindent}{0pt}
\setlength{\parskip}{3pt}
\pagestyle{empty}
\makeatletter
% figure registry (figdefs.tex is generated by the extract step)
\newcommand{\qfigdef}[4]{\@namedef{qf@file@#1}{#2}\@namedef{qf@w@#1}{#3}\@namedef{qf@wide@#1}{#4}}
\newcommand{\qf@check}[1]{\@ifundefined{qf@file@#1}{\PackageError{score2}{Unknown figure id #1}{}}{}}
\newcommand{\qf@img}[4][]{\includegraphics[width=#2\dimexpr\@nameuse{qf@w@#4}pt\relax,max width=#3,max height=0.45\textheight,#1]{\@nameuse{qf@file@#4}}}
% Full-width questions: finalize wraps every question that contains a wide figure in \qfullbegin...\qfullend,
% so the WHOLE question (stem, figures, options) is set across the page in reading order, then the two columns
% resume. The question's head - its number, stem and first figure (\qfullsplit marks the end of the head) -
% is kept on one page, so a number is never left at a page foot with its figure overleaf; the rest may break
% across pages normally, so no big gaps are left. Without a \qfullsplit, the whole question is kept together
% if it is under 45% of a page. Never write these yourself.
\newif\ifqf@full
\newif\ifqf@boxed
\newsavebox{\qf@box}
\newcommand{\qf@open}{\global\qf@fulltrue\par\global\qf@boxedtrue
  \setbox\qf@box\vbox\bgroup\hsize\textwidth\linewidth\textwidth}
\newcommand{\qf@emit}[1]{\par\egroup\global\qf@boxedfalse
  \ifdim\dimexpr\ht\qf@box+\dp\qf@box\relax<#1\textheight\noindent\usebox{\qf@box}\par\else\unvbox\qf@box\fi}
\newcommand{\qfullsplit}{\ifqf@boxed\qf@emit{0.85}\fi}
\newcommand{\qf@close}{\ifqf@boxed\qf@emit{0.45}\fi\global\qf@fullfalse}
\newcommand{\qfullbegin}{\end{multicols}\qf@open}
\newcommand{\qfullend}{\qf@close\begin{multicols}{2}}
\newcommand{\qfullstart}{\qf@open}
\newcommand{\qfulldone}{\qf@close}
% \qfig[scale]{id}: block figure, centred, directly where it is written (size set by the tool).
% Fallback if a wide figure is not inside a full-width question: break out of the columns in place - no floats.
\newcommand{\qfig}[2][1]{\qf@check{#2}%
  \ifqf@full\qf@blk{#1}{#2}\else
  \ifnum\@nameuse{qf@wide@#2}=1 \expandafter\@firstoftwo\else\expandafter\@secondoftwo\fi
  {\end{multicols}\begin{center}\qf@img{#1}{\textwidth}{#2}\end{center}\begin{multicols}{2}}%
  {\qf@blk{#1}{#2}}\fi}
\newcommand{\qf@blk}[2]{\par\smallskip{\centering\qf@img{#1}{\linewidth}{#2}\par}\smallskip}
% \qfigi[scale]{id}: inline figure (inside option lists / tables / minipages); never breaks columns.
\newcommand{\qfigi}[2][1]{\qf@check{#2}\qf@img[valign=c]{#1}{\linewidth}{#2}}
\makeatother
% structure macros - lint reads these to check completeness and figure placement
\newcommand{\papersec}[1]{\par\bigskip{\centering\Large\bfseries PAPER-#1\par}\medskip}
\newcommand{\partsec}[2]{\par\medskip\noindent{\large\bfseries PART-#1: #2}\par\smallskip}
\newcommand{\sect}[1]{\par\smallskip\noindent\textbf{#1}\par}
\newcommand{\qn}[1]{\par\smallskip\noindent\textbf{#1.}~}
\newcommand{\paraq}[2]{\par\smallskip\noindent\textbf{Paragraph for Questions #1 and #2}\par}
\newenvironment{qoptions}{\begin{itemize}[leftmargin=1.8em,itemsep=1pt,topsep=2pt,label={}]}{\end{itemize}}
\newcommand{\opt}[2]{\item[(#1)] #2}
% figure registry (generated by score2_tools.py finalize)
\qfigdef{p1_phy_s1_q01_1}{Score-II - Test 28 (23 Mar 2024)-img/p1_phy_s1_q01_1.png}{108.3}{0}
\qfigdef{p1_phy_s1_q04_1}{Score-II - Test 28 (23 Mar 2024)-img/p1_phy_s1_q04_1.png}{212.8}{0}
\qfigdef{p1_phy_s1_q05_1}{Score-II - Test 28 (23 Mar 2024)-img/p1_phy_s1_q05_1.png}{105.8}{0}
\qfigdef{p1_phy_s2_pre_1}{Score-II - Test 28 (23 Mar 2024)-img/p1_phy_s2_pre_1.png}{111.2}{0}
\qfigdef{p1_phy_s2_pre_2}{Score-II - Test 28 (23 Mar 2024)-img/p1_phy_s2_pre_2.png}{114.8}{0}
\qfigdef{p1_phy_s2_pre_3}{Score-II - Test 28 (23 Mar 2024)-img/p1_phy_s2_pre_3.png}{110.2}{0}
\qfigdef{p1_phy_s2_pre_4}{Score-II - Test 28 (23 Mar 2024)-img/p1_phy_s2_pre_4.png}{111.1}{0}
\qfigdef{p1_phy_s2_q08_1}{Score-II - Test 28 (23 Mar 2024)-img/p1_phy_s2_q08_1.png}{104.9}{0}
\qfigdef{p1_phy_s2_q09_1}{Score-II - Test 28 (23 Mar 2024)-img/p1_phy_s2_q09_1.png}{169.0}{0}
\qfigdef{p1_phy_s2_q10_1}{Score-II - Test 28 (23 Mar 2024)-img/p1_phy_s2_q10_1.png}{198.6}{0}
\qfigdef{p1_phy_s3_q01_1}{Score-II - Test 28 (23 Mar 2024)-img/p1_phy_s3_q01_1.png}{178.5}{0}
\qfigdef{p1_phy_s3_q04_1}{Score-II - Test 28 (23 Mar 2024)-img/p1_phy_s3_q04_1.png}{153.8}{0}
\qfigdef{p1_phy_s3_q05_1}{Score-II - Test 28 (23 Mar 2024)-img/p1_phy_s3_q05_1.png}{130.0}{0}
\qfigdef{p1_phy_s3_q07_1}{Score-II - Test 28 (23 Mar 2024)-img/p1_phy_s3_q07_1.png}{130.8}{0}
\qfigdef{p1_phy_s3_q08_1}{Score-II - Test 28 (23 Mar 2024)-img/p1_phy_s3_q08_1.png}{128.1}{0}
\qfigdef{p1_chem_s1_q03_1}{Score-II - Test 28 (23 Mar 2024)-img/p1_chem_s1_q03_1.png}{107.0}{0}
\qfigdef{p1_chem_s1_q04_1}{Score-II - Test 28 (23 Mar 2024)-img/p1_chem_s1_q04_1.png}{224.8}{0}
\qfigdef{p1_chem_s3_q05_1}{Score-II - Test 28 (23 Mar 2024)-img/p1_chem_s3_q05_1.png}{197.4}{0}
\qfigdef{p1_math_s3_q02_1}{Score-II - Test 28 (23 Mar 2024)-img/p1_math_s3_q02_1.png}{141.1}{0}
% ---- body skeleton ----
% \begin{document}
% \begin{center}{\Large\bfseries JEE (Main+Advanced) Score-II --- Test NN (DD Mon YYYY)}\end{center}
% \begin{multicols}{2}
%   \papersec{1}   (only for two-paper tests)
%   \partsec{1}{PHYSICS} \sect{Section-I (i): Single Correct} \qn{1} ... \qfig{p1_phy_q01_1} ...
% \end{multicols}
% \end{document}

\begin{document}
\begin{center}{\Large\bfseries JEE (Main+Advanced) Score-II --- Test 28 (23 Mar 2024)}\end{center}
\begin{multicols}{2}

\partsec{1}{PHYSICS}
\sect{Section-I(i)}

\qn{1} A photoelectric tube is subjected to an accelerating potential V as shown. The reading of ammeter is I. In which of the following cases will the ammeter reading decreases?
\qfig{p1_phy_s1_q01_1}
\begin{qoptions}
\opt{A}{An uniform magnetic field is applied parallel to the tube.}
\opt{B}{A uniform magnetic field is applied perpendicular to the plane of paper.}
\opt{C}{A uniform electric field is switched on whose direction is from emitter to collector.}
\opt{D}{A uniform electric field perpendicular to plane is applied.}
\end{qoptions}

\qn{2} A shell of radius 1m is coated with thin layer of $\beta^-$ active material. Its initial charge is zero and initial number of active atoms is $\dfrac43\times10^{12}$. If half life of decay is 1 hr and all the electrons are emitted with energy 1.44 keV. Select the correct options and neglect the time taken by e$^-$ to return back
\begin{qoptions}
\opt{A}{Maximum number of electrons emitted is $10^{10}$}
\opt{B}{Maximum number of electrons emitted is $10^{12}$}
\opt{C}{The time after which charge on shell becomes constant is 2 half lifes.}
\opt{D}{The time after which charge on shell becomes constant is 4 half lifes.}
\end{qoptions}

\qn{3} In an experiment to find focal length of an equiconvex lens, the mount readings for lens upright, object pin upright and image pin upright are respectively 50.0 cm, 31.3 cm and 80.7 cm on the meter scale. When an index needle of length 10.0 cm (measured by meter scale) is placed between lens and object pin's tip, readings of lens mount and object pin mount are 50.0 cm and 41.0 cm respectively. When the same needle is placed between lens and image pin's tip the respective mounts read 50.0 cm and 61.0 cm. Thickness of the lens is given, t = 0.60 cm measured using a vernier calliper of least count 0.01 cm. Choose the \textbf{CORRECT} options.
\begin{qoptions}
\opt{A}{Index correction for object distance is $(1.3\pm0.3)$ cm.}
\opt{B}{Index correction for image distance is $(-0.7\pm0.3)$ cm.}
\opt{C}{Focal length of lens is $(12.0\pm0.6)$ cm.}
\opt{D}{Focal length of lens is $(12.0\pm0.3)$ cm.}
\end{qoptions}

\qn{4} An unpolarised light of intensity `I$_0$' passing through two defective polarisers is shown in figure. I is the intensity of transmitted light from the second polariser. If $\alpha$ and $\beta$ be the fraction of transmitted energy of $\vec E-$vibration passing through parallel and perpendicular to transmission axis respectively, then [$\theta$ = angle between TA of both polarisers]
\qfig{p1_phy_s1_q04_1}
\begin{qoptions}
\opt{A}{For $\beta=0$, $\alpha=1$ \& $\theta=45^\circ$, $I=\dfrac{I_0}{8}$}
\opt{B}{For $\beta=\dfrac12$, $\alpha=1$, $\theta=45^\circ$, $I=\dfrac{9I_0}{16}$}
\opt{C}{For $\beta=\dfrac12$, $\alpha=\dfrac34$, $\theta=60^\circ$, $I=\dfrac{49I_0}{128}$}
\opt{D}{For $\beta=\dfrac12$, $\alpha=1$, $\theta=45^\circ$, $I=\dfrac{9I_0}{32}$}
\end{qoptions}

\qn{5} As shown in figure, two large, parallel, vertical conducting plates separated by distance `d' are charged to their potential $+$V$_0$ and $-$V$_0$. A small conducting ball of mass m and radius R(R$<<$d) hangs midway between the plates. The thread of length `L' supporting the ball is conducting wire, connected to earth, so the potential of ball is fixed at V = 0. The ball hangs straight down in equilibrium. (Given $\dfrac{1}{4\pi\varepsilon_0}=k$].
\qfig{p1_phy_s1_q05_1}
\begin{qoptions}
\opt{A}{For small horizontal displacement (x) of ball (x $<<$ L) along electric field, modulus of charge on ball is $|Q|=\dfrac{V_0xR}{kd}$}
\opt{B}{For small horizontal displacement (x) of ball (x $<<$ L) along electric field, modulus of charge on ball is $|Q|=\dfrac{2V_0xR}{kd}$}
\opt{C}{equilibrium of ball is stable if $V_0<\sqrt{\dfrac{kd^2mg}{4RL}}$}
\opt{D}{equilibrium of ball is stable if $V_0<\sqrt{\dfrac{kd^2mg}{2RL}}$}
\end{qoptions}

\qn{6} A particle of charge 2C and mass $10^{-2}$ kg is projected from origin along positive x-direction, in a magnetic field $\vec B=(-4x\hat k)$ tesla, with speed 400 m/s. When velocity becomes parallel to y-axis
\begin{qoptions}
\opt{A}{x-coordinate will be 1m}
\opt{B}{force experience by the charge will be 3200N}
\opt{C}{speed at this instant will be 600 m/s}
\opt{D}{radius of curvature of path will be 0.5m}
\end{qoptions}

\sect{Section-I(ii)}

\qn{7}\par{\small\setlength{\tabcolsep}{3pt}\noindent\begin{tabular}{|c|p{2.75cm}|c|p{1.95cm}|}
\hline
 & \centering List-I & & \centering List-II \tabularnewline
\hline
(P) & Graphical representation of pressure variation in both end open organ pipe.\newline \qfigi{p1_phy_s2_pre_1} & (1) & $\lambda=0.90$m \\
\hline
(Q) & Graphical representation of pressure variation in one end close organ pipe.\newline \qfigi{p1_phy_s2_pre_2} & (2) & $\lambda=0.96$m \\
\hline
(R) & Graphical representation of pressure variation in both end open organ pipe\newline \qfigi{p1_phy_s2_pre_3} & (3) & $\lambda=1.8$m \\
\hline
(S) & Graphical representation of pressure variation in one end open organ pipe.\newline \qfigi{p1_phy_s2_pre_4} & (4) & $\lambda=2.4$m \\
\hline
 & & (5) & $\lambda=2$m \\
\hline
\end{tabular}\par}
\begin{qoptions}
\opt{A}{P $\to$ 1;Q $\to$ 2;R $\to$ 3;S $\to$ 5}
\opt{B}{P $\to$ 4;Q $\to$ 1;R $\to$ 5;S $\to$ 2}
\opt{C}{P $\to$ 2;Q $\to$ 1;R $\to$ 5;S $\to$ 3}
\opt{D}{P $\to$ 1;Q $\to$ 5;R $\to$ 4;S $\to$ 2}
\end{qoptions}

\qn{8} A long cylinder of radius R having a cylindrical cavity of radius R/2 is shown in figure. Except cavity, the entire space inside the cylinder has a uniform magnetic field parallel to the axis of the cylinder. The magnetic field starts changing at a uniform rate of $\dfrac{dB}{dt}=k$ T/s. Point P is at a distance of R/2 from line OO$_1$.\\
(k is positive constant)
\qfig{p1_phy_s2_q08_1}
\begin{center}\small\setlength{\tabcolsep}{3pt}
\begin{tabular}{|c|p{2.6cm}|c|p{2.2cm}|}
\hline
 & \centering List-I & & \centering List-II \tabularnewline
\hline
(P) & Direction of induced electric field at Q & (1) & Parallel to line OO$_1$ \\
\hline
(Q) & Direction of induced electric field at P & (2) & $\dfrac{Rk}{3}$ \\
\hline
(R) & Magnitude of induced electric field at P & (3) & Perpendicular to line OO$_1$ \\
\hline
(S) & Magnitude of induced electric field at Q & (4) & $\dfrac{kR}{4}$ \\
\hline
 & & (5) & $\dfrac{Rk}{8}$ \\
\hline
\end{tabular}
\end{center}
\begin{qoptions}
\opt{A}{P $\to$ 3;Q $\to$ 3;R $\to$ 4;S $\to$ 5}
\opt{B}{P $\to$ 3;Q $\to$ 1;R $\to$ 2;S $\to$ 4}
\opt{C}{P $\to$ 1;Q $\to$ 3;R $\to$ 2;S $\to$ 4}
\opt{D}{P $\to$ 1;Q $\to$ 3;R $\to$ 4;S $\to$ 5}
\end{qoptions}

\qn{9} Rod of mass = 1 kg, length $\ell$ = 2m is moving with linear speed v$_0$ = 12 m/s and angular speed $\omega_0$ = 4 rad/sec about center of mass of rod on smooth horizontal surface (as shown in figure). Its one end collides with another identical rod and stick to it and start acting as single rod of length 2$\ell$. [$\pi$ = 3.14]
\qfig{p1_phy_s2_q09_1}
\begin{center}\small\setlength{\tabcolsep}{3pt}
\begin{tabular}{|c|p{3.3cm}|c|p{0.9cm}|}
\hline
 & \centering List-I & & \centering List-II \tabularnewline
\hline
(P) & Angular speed of rod just after collision (in rad/s) & (1) & 1.0 \\
\hline
(Q) & Loss in kinetic energy during collision (in joule) & (2) & 8.0 \\
\hline
(R) & Impulse exerted by rod-1 on rod-2 (in N-S) & (3) & 6.0 \\
\hline
(S) & Time taken (in sec) by combined rod to complete four rotations & (4) & 5.0 \\
\hline
 & & (5) & 5.3 \\
\hline
\end{tabular}
\end{center}
\begin{qoptions}
\opt{A}{P $\to$ 3;Q $\to$ 5;R $\to$ 1;S $\to$ 4}
\opt{B}{P $\to$ 1;Q $\to$ 3;R $\to$ 2;S $\to$ 2}
\opt{C}{P $\to$ 4;Q $\to$ 5;R $\to$ 1;S $\to$ 4}
\opt{D}{P $\to$ 4;Q $\to$ 3;R $\to$ 3;S $\to$ 4}
\end{qoptions}

\qn{10} The right limb is long enough and its radius is twice that of left limb. Water is slowly poured into the right limb. Considering water wets completely and the left limb is fully filled. List I corresponds to height of liquid in right limb and list II corresponds to radius of meniscus (R) and its shape in left limb.
\begin{center}\small\setlength{\tabcolsep}{3pt}
\begin{tabular}{|c|p{2.2cm}|c|p{2.5cm}|}
\hline
 & \centering List-I & & \centering List-II \tabularnewline
\hline
(P) & h = 1.1 cm & (1) & Concave upward \\
\hline
(Q) & h = 2 cm & (2) & Flat \\
\hline
(R) & h = 2.54 cm & (3) & Convex upward \\
\hline
(S) & h = 6.86 cm & (4) & R = 1mm \\
\hline
 & & (5) & R = 0.5 mm \\
\hline
\end{tabular}
\end{center}
\qfig{p1_phy_s2_q10_1}
\begin{qoptions}
\opt{A}{P $\to$ 3,5;Q $\to$ 2;R $\to$ 2,5;S $\to$ 1,5}
\opt{B}{P $\to$ 1,5;Q $\to$ 1;R $\to$ 2,5;S $\to$ 1,5}
\opt{C}{P $\to$ 1,5;Q $\to$ 1;R $\to$ 1,4;S $\to$ 3,5}
\opt{D}{P $\to$ 3,5;Q $\to$ 2;R $\to$ 2,4;S $\to$ 3,5}
\end{qoptions}

\sect{Section-II}

\qn{1} The figure shows a thin equiconvex lens equally divided into two pieces and arranged as shown. The refractive index of lens is 1.5 and its focal length in water is 16 cm. A plane mirror has been placed as shown 3 mm below the principal axis of the two lenses. An extended object AB of height 2mm has been placed in water 24 cm away from the left lens piece. On right side of right lens piece there is air. Find the y-coordinate (in cm) of the final image of the top end A of the object with respect to the given axis system.
\qfig{p1_phy_s3_q01_1}

\qn{2} A closed container with adiabatic walls containing an ideal monoatomic gas is being moved with constant velocity v horizontally. If the speed of the container is suddenly increased to $\dfrac{6v}{5}$ and then maintained at $\dfrac{6v}{5}$ in the same direction then the change in temperature of the gas is $\Delta T_1$ while if the speed of the container is suddenly decreased to $\dfrac{3v}{4}$ and then maintained at $\dfrac{3v}{4}$ in the same direction then the change in temperature of the gas is $\Delta T_2$. Find $\dfrac{\Delta T_1}{\Delta T_2}$

\qn{3} S is the source of Nuclear energy kept at centre of a sphere of Radius R. In steady state surface temperature of the sphere is maintained at T$_0$ when the sphere and source system is alone. Now this sphere is surrounded by a perfectly conducting spherical shell of radius 2R. Assume all surfaces are black bodies. (Assume outside of shell is vacuum having zero temperature). If the surface temperature of sphere of Radius R in equilibrium is $\left(\dfrac xy\right)^{\frac14}T_0$. Find value of (x + y). x \& y are smallest integer.

\qn{4} $\varepsilon_0$ = 20.5 V \quad $\varepsilon_1$ = 26V \quad r$_1$ = 1$\Omega$\\
r = 0.1$\Omega$ \quad\quad $\varepsilon_2$ = 2V\\
In a potentiometer setup a battery $\varepsilon_0$ is connected across a 2 metre long uniform wire having resistance 2 ohm/meter. If a null point (N) is obtained when AN = 50 cm then find value of resistance r$_2$ (in ohms)
\qfig{p1_phy_s3_q04_1}

\qn{5} Water is flowing through a hose pipe as shown in figure. If volume flow rate of water is 0.025 m$^3$/s, find force (in KN) offered by clamp to hold the pipe. Assume that pressure at cross section 2 is 0 pascal.
\qfig{p1_phy_s3_q05_1}

\qn{6} A non-rigid dielectric of dielectric constant K = 5, Young's module $4.43\times10^3\ \dfrac{\text N}{\text m^2}$ and thickness 10 cm is maintain the separation between the two plates in a parallel plate capacitor. If applied potential difference is nearly zero than capacitance of capacitor is C$_0$ but at higher potential the plates get compressed by small amount due to attractive force between the plates. The fractional change in capacitance is N $\times10^{-7}$ when a battery of 1 kV is connected across it. Find the value of N. (assume change in K is negligible) $\left\{\varepsilon_0=8.86\times10^{-12}\ \dfrac{\text C^2}{\text{N--m}^2}\right\}$

\qn{7} A wedge of mass 4m is placed on a block of same mass which is placed on a rough horizontal surface with coefficient of friction $\mu_2$ = 0.2. The coefficient of friction between the block and the wedge is $\mu_1$ = 0.4. A ball of mass m moving with velocity v$_0$ = 10 m/s collides with the wedge as shown and sticks to it. Find the sum of velocities of the wedge and the block just after collision in m/s. Given that the wedge does not topple.
\qfig{p1_phy_s3_q07_1}

\qn{8} A straight horizontal conducting string of length $\ell$ has one end fixed to a vertical fixed cylinder of radius R at point A as shown. On other end a particle is attached. As shown in the top view there is a uniform downward magnetic field B in the whole region. The particle is projected perpendicular to the string with speed v along the horizontal plane. Friction coefficient between particle and floor is $\mu$. Find the average emf developed (in volt) in the string during its motion till it collides with the cylinder.\\
Given $\mu$ = 0.2, $\ell$ = 1m, R = 2m, B = 3T, v = 1 m/s
\qfig{p1_phy_s3_q08_1}

\partsec{2}{CHEMISTRY}
\sect{Section-I(i)}

\qn{1} If electron of hydrogen atom is replaced by another particle of same charge but of double mass then : [Assume that nucleus is stationary]
\begin{qoptions}
\opt{A}{Bohr's radius will increase.}
\opt{B}{Ionisation energy of atom will be doubled.}
\opt{C}{Speed of new particle in given state will be lesser than electrons speed in same orbit.}
\opt{D}{Gap between energy levels will now be doubled.}
\end{qoptions}

\qn{2} Select the correct statement(s) :
\begin{qoptions}
\opt{A}{Primitive arrangement is observed in all seven crystal systems.}
\opt{B}{Carborundum is hard and brittle covalent solid.}
\opt{C}{In CrO$_2$ magnetic moment of the domains are aligned in parallel and anti parallel direction in unequal numbers.}
\opt{D}{In crystalline solids density of defects per unit volume increases on increase in temperature.}
\end{qoptions}

\qn{3} Which statement(s) is/are CORRECT?
\qfig{p1_chem_s1_q03_1}
\begin{qoptions}
\opt{A}{Reduced product of P and Q will be metamers to each other. (Consider LiAlH$_4$ as a reducing agent)}
\opt{B}{By dry distillation of hydrolysed product of P with Ca(OH)$_2$ gives benzophenone.}
\opt{C}{One of the hydrolysis product of Q reacts with NaNO$_2$ + HCl followed by reaction with phenol in mild basic medium gives orange red dye.}
\opt{D}{Electrophile involved in the formation of Q is dichlorocarbene.}
\end{qoptions}

\Needspace{9\baselineskip}\qn{4} Considering the following reaction sequence, the correct statement(s) is(are)
\qfig{p1_chem_s1_q04_1}
\begin{qoptions}
\opt{A}{Compounds \textbf{P} and \textbf{Q} are carboxylic acids.}
\opt{B}{Compound \textbf{S} decolorizes bromine water.}
\opt{C}{Compounds \textbf{P} and \textbf{S} react with hydroxylamine to give the corresponding oximes.}
\opt{D}{Compound \textbf{R} reacts with dialkylcadmium followed by NaBH$_4$, MeOH gives secondary alcohol.}
\end{qoptions}

\qn{5} Which of the following statement(s) is/are correct
\begin{qoptions}
\opt{A}{Due to (s--p) intermixing in N$_2$--molecule, energy of $\sigma_{2p}$ increases}
\opt{B}{If a molecular orbital is symmetrical around the bond axis, can be designated as $\sigma$*M.O.}
\opt{C}{Both bond energy and bond length is higher for O$_2$ as compare to C$_2$ molecule}
\opt{D}{Both B$_2$ and O$_2$ molecules contain same number of unpaired electrons.}
\end{qoptions}

\qn{6} In which of the following reactions PbO$_2$ will not form
\begin{qoptions}
\opt{A}{PbS + boiled and conc. HNO$_3\to$}
\opt{B}{Pb$^{+2}\ \xrightarrow[\text{(ii) H}_2\text{O}_2]{\text{(i) excess NaOH}}$}
\opt{C}{Pb$^{+2}$ + CrO$_4^{2-}\to$}
\opt{D}{Pb + conc. HNO$_3\to$}
\end{qoptions}

\sect{Section-I(ii)}

\qn{7} Electrolysis of following aqueous solution proceed in electrolytic cell for small period of time.
\begin{center}\small\setlength{\tabcolsep}{3pt}
\begin{tabular}{|c|p{2.5cm}|c|p{2.5cm}|}
\hline
\multicolumn{2}{|c|}{\parbox{3.3cm}{\centering List-I\\(Aqueous solution)}} & \multicolumn{2}{c|}{\parbox{3.3cm}{\centering List-II\\(Products)}} \\
\hline
(P) & CH$_3$COOK (Pt electrodes) & (1) & Metal becomes deposited at cathode \\
\hline
(Q) & Dilute NaCl (Pt electrodes) & (2) & H$_2$ gas at cathode \\
\hline
(R) & CuSO$_4$ (Pt electrodes) & (3) & H$_2$ gas at anode \\
\hline
(S) & Ni(NO$_3$)$_2$ (Ni electrodes) & (4) & O$_2$ gas at anode \\
\hline
 & & (5) & Concentration of electrolyte remain almost same \\
\hline
\end{tabular}
\end{center}
\begin{qoptions}
\opt{A}{P $\to$ 1,4;Q $\to$ 2,4,5;R $\to$ 3,4;S $\to$ 1,5}
\opt{B}{P $\to$ 2,4;Q $\to$ 2,4,5;R $\to$ 3,4;S $\to$ 1,5}
\opt{C}{P $\to$ 2;Q $\to$ 2,4,5;R $\to$ 1,4;S $\to$ 1,5}
\opt{D}{P $\to$ 2;Q $\to$ 1,4,5;R $\to$ 1,5;S $\to$ 1,4}
\end{qoptions}

\Needspace{12\baselineskip}\qn{8}\par\nopagebreak[4]
\begin{center}\small\setlength{\tabcolsep}{3pt}
\begin{tabular}{|c|p{2.7cm}|c|p{2.0cm}|}
\hline
 & \centering List-I & & \centering List-II \tabularnewline
\hline
(P) & Process of precipitating colloidal sol by an electrolyte & (1) & Peptization \\
\hline
(Q) & Aggregation of ions & (2) & Coagulation \\
\hline
(R) & Process of formation colloidal sol by an electrolyte from precipitate & (3) & Micellization \\
\hline
(S) & Process of removal of dissolved impurity from colloidal solution & (4) & Dialysis \\
\hline
 & & (5) & Electro-osmosis \\
\hline
\end{tabular}
\end{center}
\begin{qoptions}
\opt{A}{P $\to$ 3;Q $\to$ 2;R $\to$ 1;S $\to$ 4}
\opt{B}{P $\to$ 2;Q $\to$ 3;R $\to$ 1;S $\to$ 4}
\opt{C}{P $\to$ 2;Q $\to$ 1;R $\to$ 4;S $\to$ 3}
\opt{D}{P $\to$ 1;Q $\to$ 3;R $\to$ 2;S $\to$ 4}
\end{qoptions}

\qn{9} Match each species in List-I with List - II and choose correct option
\begin{center}\small\setlength{\tabcolsep}{3pt}
\begin{tabular}{|c|p{3.0cm}|c|p{1.4cm}|}
\hline
\multicolumn{2}{|c|}{List-I} & \multicolumn{2}{c|}{List-II} \\
\hline
(I) & A carbohydrate which yields only glucose on hydrolysis & (P) & Sucrose \\
\hline
(II) & A carbohydrate which yields glucose and fructose on hydrolysis & (Q) & Lactose \\
\hline
(III) & A carbohydrate which yields glucose and galactose on hydrolysis & (R) & Amylose \\
\hline
(IV) & A carbohydrate which can reduce Fehling's solution & (S) & Mannose \\
\hline
 & & (T) & Cellulose \\
\hline
\end{tabular}
\end{center}
\begin{qoptions}
\opt{A}{I $\to$ R,T;II $\to$ Q;III $\to$ P;IV $\to$ Q,S}
\opt{B}{I $\to$ T;II $\to$ Q;III $\to$ Q;IV $\to$ Q}
\opt{C}{I $\to$ R;II $\to$ P;III $\to$ Q;IV $\to$ P}
\opt{D}{I $\to$ R,T;II $\to$ P;III $\to$ Q;IV $\to$ Q,S}
\end{qoptions}

\qn{10}\par\nopagebreak[4]
\begin{center}\small\setlength{\tabcolsep}{3pt}
\begin{tabular}{|c|p{2.2cm}|c|p{2.9cm}|}
\hline
 & \centering List-I (Property of any product) & & \centering List-II (Reaction) \tabularnewline
\hline
(P) & Di-atomic gas released in reaction & (1) & B$_2$H$_6$ + NH$_3$ (1 : 2) $\xrightarrow{\Delta}$ \\
\hline
(Q) & Volatile compound burn with green edged flame & (2) & B$_2$H$_6$ + MeOH $\to$ \\
\hline
(R) & Aromatic compound will form in reaction & (3) & BCl$_3$ + NH$_4$Cl $\xrightarrow[\text{PhCl}]{\text{420K}}$ X $\xrightarrow{\text{NaBH}_4}$ \\
\hline
(S) & Back bonding is present in boron compound & (4) & NaBH$_4$ + I$_2$ $\to$ \\
\hline
 & & (5) & H$_3$BO$_3$ + MeOH $\to$ \\
\hline
\end{tabular}
\end{center}
The the most appropriate match.
\begin{qoptions}
\opt{A}{P $\to$ 1,4;Q $\to$ 5;R $\to$ 1,3;S $\to$ 2,5}
\opt{B}{P $\to$ 4,5;Q $\to$ 2;R $\to$ 1;S $\to$ 1,3}
\opt{C}{P $\to$ 1;Q $\to$ 4,5;R $\to$ 3,5;S $\to$ 1,3,5}
\opt{D}{P $\to$ 4;Q $\to$ 2,5;R $\to$ 2,3;S $\to$ 1,2}
\end{qoptions}

\sect{Section-II}

\qn{1} 250 mL of 0.16 M MgCl$_2$(aq) is mixed with 750 mL of 0.2 M NH$_3$(aq) resulting in precipitation of some Mg(OH)$_2$. To redissolve Mg(OH)$_2$, what minimum mass (in gram) of (NH$_4$)$_2$SO$_4$ should be added to the above solution?\\
(Given : K$_{sp}$ [Mg(OH)$_2$] = $9\times10^{-12}$, K$_b$ (NH$_3$) = $2\times10^{-5}$)

\qn{2} For reaction A(g) + B(g) $\rightleftharpoons$ C(g), at 300 K, the rate of backward reaction at time, t, is 10 times of forward reaction for the reaction. Calculate free energy change (in cal) involved at that time of reaction.\\
(Given : lnx = 2.3 logx)

\qn{3} [A] is a monomer of natural rubber, on reaction with HBr (1 eq)/CCl$_4$ at 298 K, it forms major product [B]. The number of allylic hydrogen of [A] is x and that of [B] is y, the value of (y $-$ x)?

\qn{4} Reaction of \textbf{x} g of Sn with HCl quantitatively produced a salt. Entire amount of the salt reacted with \textbf{y} g of nitrobenzene in the presence of required amount of HCl to produce 2.58 g of an organic salt (quantitatively).\\
(Use Molar masses (in g mol$^{-1}$) of H, C, N, O, Cl and Sn as 1, 12, 14, 16, 35 and 119, respectively).\\
The value of (x + y) is :

\qn{5} The number of sp$^2$ carbons in major product formed from the following reaction sequence is \underline{\hspace{1.2cm}}.
\qfig{p1_chem_s3_q05_1}

\qn{6} Cl$_{2(g)}$ reacts with excess of F$_{2(g)}$ at 573 K to form interhalogen ClF$_{x(g)}$. The value of x is ...........

\qn{7} When three parts of concentrated HCl and one part of concentrated HNO$_3$ are mixed, aqua regia is formed which is used for dissolving noble metals. If 3 mole Pt metal are dissolved in aqua regia then `X' mole of colorless gas is produced and if gold is dissolved in aqua regia then it forms a complex in `Y' oxidation state of Au. Calculate the sum of `X' and `Y'.

\qn{8} Calculate sum of unpaired electrons in per molecule of all the paramagnetic products for following reactions.\\
(i) Pyrolusite + KOH $\xrightarrow{\text{O}_2}$\\
(ii) Pb(NO$_3$)$_2$ $\xrightarrow{\text{673K}}$\\
(iii) Ni$^{2+}_{\text{(aq.)}}$ + dmg$^-$ $\to$

\partsec{3}{MATHEMATICS}
\sect{Section-I(i)}

\qn{1} Let A$\left(\dfrac12,1,3\right)$ and C$\left(3,\dfrac{17}{4},\dfrac{13}{2}\right)$ be the vertices of $\Delta$ABC and equation of angle bisector of $\angle$ABC is $\dfrac{x-1}{1}=\dfrac{y-2}{1}=\dfrac{z-3}{2}$, then
\begin{qoptions}
\opt{A}{equation of side BC is $\dfrac{x-1}{4}=\dfrac{2y-5}{7}=\dfrac{z-2}{9}$}
\opt{B}{equation of side BC is $\dfrac{x-1}{2}=\dfrac{y+2}{7}=\dfrac{z+1}{3}$}
\opt{C}{co-ordinate of vertex B is $(1,-2,-1)$}
\opt{D}{co-ordinate of vertex B is $(5,6,11)$}
\end{qoptions}

\qn{2} Let the minimum value of $\sqrt{\alpha^2+(\alpha+1)^2}+\sqrt{(\alpha-1)^2+(\alpha-2)^2}$ be $\lambda$ and the equation of the two sides of a triangle are $4x+5y-20=0$ and $x-y+4=0$. If the circum-centre of the triangle be P$\left(\sqrt{\dfrac52}\lambda,0\right)$, then
\begin{qoptions}
\opt{A}{centroid of the triangle is $\left(\dfrac{11}{3},\dfrac53\right)$}
\opt{B}{orthocentre of the triangle is (1, 5)}
\opt{C}{area of the triangle is 9 sq. units}
\opt{D}{circum-radius of the triangle is $\sqrt{41}$ units}
\end{qoptions}

\qn{3} If the curve represented by differential equation $(y^4+1)dx+4y^3dy=e^{2y-x}dy$; passes through origin then which of the following statement(s) is(are) true ?
\begin{qoptions}
\opt{A}{$\displaystyle\lim_{y\to\infty}\dfrac{e^x}{e^{2y}}=0$}
\opt{B}{The line $y=x$ meets the curve at three distinct points}
\opt{C}{The curve cuts the y-axis only at origin}
\opt{D}{The line $y=2x+3$ touches the curve at origin}
\end{qoptions}

\qn{4} Consider a matrix $M=\begin{bmatrix}x&y&z\\x^2&y^2&z^2\\y+z&z+x&x+y\end{bmatrix}$ if $\dfrac{\det(\mathrm{adj}(\mathrm{adj}(\mathrm{adj}(\mathrm{adj}\,M))))}{(x-y)^{16}(y-z)^{16}(z-x)^{16}}=2^{32}\cdot3^{16}\cdot5^{16}$ then number of such 3-triples of (x, y, z) is(are) (where $x\ne y\ne z$)
\begin{qoptions}
\opt{A}{$^{59}C_2-88$ when x, y, z $\in$ Natural number}
\opt{B}{$^{59}C_2-85$ when x, y, z $\in$ Natural number}
\opt{C}{$^{62}C_2-94$ when x, y, z $\in$ Whole number}
\opt{D}{$^{62}C_2-91$ when x, y, z $\in$ whole number}
\end{qoptions}

\qn{5} Let A \& B are two events such that $P(A\cap B)=\dfrac14$ \& $P(B)=\dfrac13$ and if $P\left(\left(A\cap B\right)/\left(A\cup\overline B\right)\right)=\dfrac{K}{22}$ then K is divisible by
\begin{qoptions}
\opt{A}{2}
\opt{B}{3}
\opt{C}{4}
\opt{D}{8}
\end{qoptions}

\qn{6} Consider the function
\[f(x)=\int_0^x\left(5\ell n(1+t^2)-10t\tan^{-1}t+16\sin t\right)dt,\]
then
\begin{qoptions}
\opt{A}{$f(x)$ is increasing $\forall\ x\in(0,1)$}
\opt{B}{$f(x)$ is decreasing $\forall\ x\in(0,1)$}
\opt{C}{$\displaystyle\int_0^xf(t)dt$ is positive $\forall\ x\in(0,1)$}
\opt{D}{$\displaystyle\int_0^xf(t)dt$ is increasing for $x\in(0,1)$}
\end{qoptions}

\sect{Section-I(ii)}

\qn{7} Let z be a complex number satisfying $z^4-6z^3+18z^2-30z+25=0$ has four roots z$_1$, z$_2$, z$_3$ and z$_4$. Where $|z_1-1-2i|=0$, Re(z$_4$) = Re(z$_1$) and Img(z$_2$) $>$ 0 and $\bar z$ denotes the complex conjugate of z. Match each entry in List-I to the correct entries in List-II.
\begin{center}\small\setlength{\tabcolsep}{3pt}
\begin{tabular}{|c|p{3.5cm}|c|p{0.6cm}|}
\hline
\multicolumn{2}{|c|}{List-I} & \multicolumn{2}{c|}{List-II} \\
\hline
(P) & The area of the quadrilateral bounded by z$_1$, z$_2$, z$_3$ \& z$_4$ is & (1) & 1 \\
\hline
(Q) & $\left|\dfrac{(z_0-\bar z_1)(z_0-\bar z_3)}{(z_0-\bar z_4)(z_0-\bar z_2)}\right|$ is equal to (where z$_0$ is any purely real no.) & (2) & 2 \\
\hline
(R) & If $(|\alpha-z_1|+\allowbreak|\alpha-z_2|+\allowbreak|\alpha-z_3|+\allowbreak|\alpha-z_4|)$ is minimum then $3|\alpha|$ is (where $\alpha$ is a complex number) & (3) & 3 \\
\hline
(S) & If $\arg\left(\dfrac{\beta-z_1}{\beta-z_3}\right)=\tan^{-1}\left(\dfrac{2}{1+\sqrt5}\right)+\tan^{-1}\left(\dfrac{1}{2+\sqrt5}\right)$ then $\dfrac{1}{\sqrt5}|\beta-2-i|_{\max}$ is & (4) & 4 \\
\hline
 & & (5) & 5 \\
\hline
\end{tabular}
\end{center}
Then the correct option is
\begin{qoptions}
\opt{A}{P $\to$ 3;Q $\to$ 3;R $\to$ 4;S $\to$ 3}
\opt{B}{P $\to$ 3;Q $\to$ 1;R $\to$ 3;S $\to$ 4}
\opt{C}{P $\to$ 3;Q $\to$ 1;R $\to$ 5;S $\to$ 2}
\opt{D}{P $\to$ 1;Q $\to$ 3;R $\to$ 2;S $\to$ 5}
\end{qoptions}

\qn{8} Match List-I with List-II. Here [$*$] denotes the greatest integer function.\\
\textbf{List-I}
\begin{description}[leftmargin=1.8em,itemsep=2pt,topsep=2pt,labelwidth=1.6em]
\item[(P)] $\displaystyle\int_{-1}^{1}\left[\dfrac13\left(\sin^{-1}x+\cos^{-1}x+\tan^{-1}x\right)\right]dx$, equals
\item[(Q)] If $f(x)=\displaystyle\lim_{k\to\infty}\int_0^{k[x]}(kt-[kt])^k dt$; $k\in N$, then the value of $f(\alpha)$, where $\alpha$ is the number of value of x satisfying the equation $2^x+3^x+4^x+5^x-7^x=0$, $x\in R$
\item[(R)] If $m=\displaystyle\int_{-10}^{0}\dfrac{\left|\frac{[x]}{3x-[x]}\right|}{\left(\frac{2[x]}{3x-[x]}\right)}dx$ then $[m]+5$ is equal to
\item[(S)] If $\alpha=\displaystyle\lim_{x\to0^-}\dfrac{[x]+[x^2]+[x^3]+\dots+[x^{2n+1}]+n+1}{1+[x]+|x|+2x}$ then
\[\resizebox{\linewidth}{!}{$\displaystyle\int_\alpha^{2\alpha^2}\sin^{-1}\!\left(\sin^{-1}\!\left(\sin^{-1}\!\left(\sin^{-1}x\right)\right)\right)dx-\left[\int_0^\pi[\cos x]\,dx\right]$}\]
equals
\end{description}
\textbf{List-II}\\
(1) 2 \quad (2) 9 \quad (3) 0 \quad (4) 7 \quad (5) 1
\begin{qoptions}
\opt{A}{P $\to$ 2;Q $\to$ 3;R $\to$ 5;S $\to$ 4}
\opt{B}{P $\to$ 3;Q $\to$ 5;R $\to$ 2;S $\to$ 1}
\opt{C}{P $\to$ 3;Q $\to$ 1;R $\to$ 2;S $\to$ 4}
\opt{D}{P $\to$ 2;Q $\to$ 5;R $\to$ 3;S $\to$ 4}
\end{qoptions}

\qn{9} Let the lines L$_1$ : $7x-y-9=0$ and L$_2$ : $x+7y+13=0$ are the equation of tangents to a parabola at the point P(2, 5) and Q($-$4/3, $-$5/3) respectively,
\begin{center}\small\setlength{\tabcolsep}{3pt}
\begin{tabular}{|c|p{3.0cm}|c|p{1.0cm}|}
\hline
 & \centering List-I & & \centering List-II \tabularnewline
\hline
(P) & If coordinate of the focus of the parabola is (h, k) then $|3h+8k|$ is & (1) & 7 \\
\hline
(Q) & If the equation of directrix of parabola is $ax+by-24=0$ then $(a-b)$ is & (2) & 9 \\
\hline
(R) & If the equation of the axis of the parabola is $\ell x+my+13=0$ then $(\ell-m)$ is & (3) & 11 \\
\hline
(S) & If tangents at the points P \& Q intersects each other at R, then the equation of circumcircle of $\Delta$PQR is $x^2+y^2+\alpha x+\beta y+\gamma=0$ then $|6\alpha+6\beta-\gamma|$ is & (4) & 13 \\
\hline
 & & (5) & 15 \\
\hline
\end{tabular}
\end{center}
\begin{qoptions}
\opt{A}{P $\to$ 1;Q $\to$ 4;R $\to$ 3;S $\to$ 4}
\opt{B}{P $\to$ 3;Q $\to$ 1;R $\to$ 2;S $\to$ 4}
\opt{C}{P $\to$ 1;Q $\to$ 2;R $\to$ 3;S $\to$ 5}
\opt{D}{P $\to$ 3;Q $\to$ 4;R $\to$ 2;S $\to$ 4}
\end{qoptions}

\qn{10} Consider the matrix $A=\begin{bmatrix}3&-4\\1&-1\end{bmatrix}$ and $B=\begin{bmatrix}a&6\\0&1\end{bmatrix}$ and Let P be any orthogonal matrix and $Q=PAP^T$ and $R=P^TQ^KP$ also $S=PBP^T$ and $T=P^TS^KP$
\begin{center}\small\setlength{\tabcolsep}{3pt}
\begin{tabular}{|c|p{2.6cm}|c|p{2.3cm}|}
\hline
 & \centering List-I & & \centering List-II \tabularnewline
\hline
(P) & If we vary K from 1 to `n', then first row first column elements of R will form & (1) & G.P. with common ratio a \\
\hline
(Q) & If we vary k from 1 to n, then 2$^{\text{nd}}$ row 2$^{\text{nd}}$ column elements of R will form & (2) & A.P with common difference 2 \\
\hline
(R) & If we vary k from 1 to n, then first row first column elements of T will form & (3) & G.P. with common ratio `6' \\
\hline
(S) & If we vary k from 3 to n, then first row 2$^{\text{nd}}$ column elements of T will represent the sum of & (4) & A.P with common difference $-2$ \\
\hline
 & & (5) & A.P with common difference $-7$ \\
\hline
\end{tabular}
\end{center}
\begin{qoptions}
\opt{A}{P $\to$ 5;Q $\to$ 3;R $\to$ 4;S $\to$ 2}
\opt{B}{P $\to$ 4;Q $\to$ 2;R $\to$ 4;S $\to$ 5}
\opt{C}{P $\to$ 2;Q $\to$ 4;R $\to$ 1;S $\to$ 1}
\opt{D}{P $\to$ 3;Q $\to$ 1;R $\to$ 2;S $\to$ 5}
\end{qoptions}

\sect{Section-II}

\qn{1} Consider the two sets of observations as follows :
\begin{center}\small\setlength{\tabcolsep}{3pt}
\begin{tabular}{|c|c|c|c|}
\hline
\textbf{Set} & \textbf{Size} & \textbf{Mean} & \textbf{Coefficient of variation} \\
\hline
D$_1$ & 13 & 10 & 20 \\
\hline
D$_2$ & 27 & 30 & 10 \\
\hline
\end{tabular}
\end{center}
If each observation of set D$_1$ and D$_2$ is multiplied by 2 and then subtracted by 10. Then the value of combined variance of new sets is

\qn{2} In the given below figure. Circles C$_1$ and C$_2$ both of radius 360 units are tangent to each other and both tangent to a straight line $\ell$. If circle C$_3$ is tangent to C$_1$, C$_2$, $\ell$ and circle C$_4$ is tangent to C$_1$, C$_3$ and $\ell$. Then radius of C$_4$ is
\qfig{p1_math_s3_q02_1}

\qn{3} Let $S_m(x)=\displaystyle\sum_{n=0}^{m}{}^mC_n\sin(nx)\cos((m-n)x)$ then the value of $\left[\dfrac{S_8(\pi/16)}{7}\right]$. (where [$*$] denotes greatest integer function)

\qn{4} The number of integral solutions (x, y) of $\tan^{-1}x+\tan^{-1}8=\tan^{-1}y+\dfrac\pi2$ is

\qn{5} If $\alpha$ is the number of points where $f(x)=\max\{|x+5|,\ |x+4|,\ |x+3|,\ |x+2|,\ |x+1|,\ |x|\}$ is non differentiable in $-6<x<1$ then number of points where $g(x)=|x^{2\alpha}-1|$ is non differentiable is

\qn{6} Let $x\in\left(0,\dfrac\pi2\right)$, $f(x)=\cot^{-1}(2-\tan x)$ and $g(\theta)=\displaystyle\int_{\min.(\theta,\,\pi/4)}^{\max.(\theta,\,\pi/4)}f(x)\,dx$ then the value of $\left(\dfrac{\pi^2}{g\left(\frac{\pi}{12}\right)-g\left(\frac\pi3\right)}\right)$ is

\qn{7} The volume of parallelopiped whose coterminous edges are represented by the vectors $3\left(\vec b\times\vec c\right)$, $2\left(\vec a\times\vec b\right)$ and $4\left(\vec c\times\vec a\right)$ where\\
$\vec a=(1+\sin\theta)\hat\imath+(\cos\theta)\hat\jmath+(\sin2\theta)\hat k$\\
$\vec b=\sin\left(\theta+\frac{2\pi}{3}\right)\hat\imath+\cos\left(\theta+\frac{2\pi}{3}\right)\hat\jmath+\sin\left(2\theta+\frac{4\pi}{3}\right)\hat k$\\
$\vec c=\sin\left(\theta-\frac{2\pi}{3}\right)\hat\imath+\cos\left(\theta-\frac{2\pi}{3}\right)\hat\jmath+\sin\left(2\theta-\frac{4\pi}{3}\right)\hat k$\\
is 18 cubic units. Then the number of value(s) of $\theta\in\left(0,\dfrac\pi2\right)$ is(are)

\qn{8} If
\[\resizebox{\linewidth}{!}{$\displaystyle\int e^{e^x+x^3+x}\left(\left(xe^x+3x^3+x+1\right)\sin x+x\cos x\right)dx=f(x)\,e^{e^x+x^3+x}+C,$}\]
(where $f\left(\dfrac\pi6\right)=\dfrac\pi{12}$ and C is the constant of integration and [*] denotes Greatest Integer Function) then $\left[f\left(\dfrac{5\pi}{2}\right)\right]$ is equal to

\end{multicols}
\end{document}
